\problemname{}

\illustration{0.3}{etymologie.jpg}{
    \emph{Van Dale Groot etymologisch woordenboek}.
    \textcopyright{}~Van~Dale~Uitgevers, used with permission
}

\newcommand{\protagonist}{Eelco}

\protagonist{} has recently started to gain interest in the field that studies the origin of words: etymology.
He especially likes how words can evolve in many different ways:
pronunciation changes over time,
words are borrowed from different languages,
and the meaning of words can change based on culture.
\protagonist{} is eager to attend the Networking With Etymologists: Revolutionary Conference for the first time ever.
To make a good first impression,
he is going to present a completely new method to make new words from existing words.

To make a new word from an existing word $s$,
\protagonist{} proposes to take every second letter of $s+s$,
starting with the first letter.
For example, applying this method to the word ``\texttt{etymology}'' would result in ``\texttt{eyooytmlg}''.
Of course, to design even more words, this process can be repeated many times.
\protagonist{} would like to prepare a list of new words to present at the conference,
so he writes a program that applies his method some predetermined number of times.

\begin{Input}
    The input consists of:
    \begin{itemize}
        \item One line with two integers $n$ and $k$ ($1 \leq n \leq 10^5$, $1 \leq k \leq 10^{18}$),
            the length of the original word and the number of times to apply the method.
        \item One line with a string $s$ of length $n$, only consisting of English lowercase letters (\texttt{a-z}),
            the original word.
    \end{itemize}
\end{Input}

\begin{Output}
    Output the resulting word after applying the method to the original word $k$ times.
\end{Output}
